% =============================================================================
%  src/cvstyle.tex  |  look & feel of the PDF (fonts + layout macros)
% -----------------------------------------------------------------------------
%  This file holds DESIGN ONLY.  Content lives in data/cv.yaml, and the colours
%  / page geometry / PDF metadata are generated into src/generated/theme.tex
%  (so the PDF palette stays identical to the web page in web/style.css).
%
%  Layer 1  src/cv.tex          -> document class, order of the sections
%  Layer 2  src/cvstyle.tex     -> this file: fonts + reusable macros
%  Layer 3  src/generated/*.tex -> AUTO-GENERATED content   (do not edit)
% =============================================================================

% ------------------------------- Fonts --------------------------------------
% Source Serif Pro ships with TeX Live, so the local build and the CI build
% produce exactly the same PDF (no system-font guessing / no silent fallback).
% The web page loads "Source Serif 4" from Google Fonts: the same typeface.
\IfFileExists{sourceserifpro.sty}{%
  \usepackage{sourceserifpro}%
}{%
  % Safety net for a minimal TeX install without the package.
  \usepackage{fontspec}%
  \IfFontExistsTF{DejaVu Serif}{\setmainfont{DejaVu Serif}}{\setmainfont{Latin Modern Roman}}%
}

% ------------------------- Packages & spacing -------------------------------
\usepackage[a4paper]{geometry}          % real values come from generated/theme.tex
\usepackage{array,tabularx}
\usepackage[table]{xcolor}
\usepackage{enumitem}
% Keeps a heading with the lines that belong to it: without this, a section bar
% or an entry headline can end up alone at the bottom of a page (the PDF then
% reads nothing like the single continuous column of the web page).
\usepackage{needspace}
\usepackage[hidelinks]{hyperref}

\setlength{\parindent}{0pt}
\setlength{\parskip}{2pt}
\renewcommand{\arraystretch}{1.05}
\tolerance=1500
\emergencystretch=1.5em
\hbadness=10000

% --------------------------- Helper primitives ------------------------------
% Flexible, ragged-right column (tabularx); R = right aligned.
\newcolumntype{Y}{>{\raggedright\arraybackslash}X}
\newcolumntype{R}{>{\raggedleft\arraybackslash}X}

% XeTeX-native letter-spacing (microtype's \textls does NOT work with xetex).
\newcommand{\cvls}[2][10]{\addfontfeatures{LetterSpace=#1}#2\addfontfeatures{LetterSpace=0}}

% Bullet used between the contact items in the header.
\newcommand{\cvsep}{\hspace{0.55em}\textcolor{cvrule}{$\bullet$}\hspace{0.55em}}

% Hairline rule (under the header).
\newcommand{\cvrule}{%
  \par\addvspace{8pt}%
  \noindent\textcolor{cvrule}{\rule{\textwidth}{0.8pt}}%
  \par\addvspace{1pt}}

% ------------------------------- Header -------------------------------------
\newcommand{\cvname}[1]{{\fontsize{25}{29}\selectfont\bfseries\cvls[14]{#1}}}
% \cvrole{Senior Software Engineer}: uppercase + letter-spaced. A letter-spaced
% font widens the gap between two WORDS exactly as much as the gap between two
% letters, so at LetterSpace=30 the line reads as one long word; \spaceskip
% gives the word gap its own width (~2.4x the letter gap, like \cvname below).
% The setting is local to this group, so nothing else in the document changes.
\newcommand{\cvrole}[1]{%
  {\large\bfseries\color{cvaccent}%
   \spaceskip=0.8em\relax
   \cvls[30]{\MakeUppercase{#1}}}}

% --------------------------- Section heading ---------------------------------
% Solid colour bar, full text width, white letter-spaced text.
\newcommand{\cvsection}[1]{%
  \par\addvspace{9pt}%
  \needspace{3\baselineskip}%  a bar never sits alone at the bottom of a page
  {\setlength{\fboxsep}{5pt}%
   \noindent\colorbox{cvbar}{%
     \makebox[\dimexpr\textwidth-2\fboxsep\relax][l]{%
       \hspace{1pt}\textcolor{cvbartext}{\cvls[9]{\textbf{\MakeUppercase{#1}}}}}}}%
  \par\addvspace{4.5pt}}

% ------------------------------- Entries -------------------------------------
% \cventry{title (bold accent)}{subtitle (italic muted)}{dates (right)}
\newcommand{\cventry}[3]{%
  \par\addvspace{4pt}%
  \needspace{3\baselineskip}%  headline + the first line of its content stay together
  \noindent
  \begin{tabularx}{\textwidth}{@{}Y r@{}}
    {\bfseries\color{cvaccent}#1} & {\footnotesize\color{cvmuted}#3}\\
    \multicolumn{2}{@{}l@{}}{\itshape\color{cvmuted}#2}
  \end{tabularx}%
  \par\addvspace{0pt}}

% Same, without a subtitle line.
\newcommand{\cventryplain}[2]{%
  \par\addvspace{4pt}\noindent
  \begin{tabularx}{\textwidth}{@{}Y r@{}}
    {\bfseries\color{cvaccent}#1} & {\footnotesize\color{cvmuted}#2}
  \end{tabularx}%
  \par\addvspace{0pt}}

% \cvproject{name}{dates} - a project headline (slightly larger than \cventry).
\newcommand{\cvproject}[2]{%
  \par\addvspace{4pt}\noindent
  \begin{tabularx}{\textwidth}{@{}Y r@{}}
    {\large\bfseries\color{cvaccent}#1} & {\footnotesize\color{cvmuted}#2}
  \end{tabularx}%
  \par\addvspace{0pt}}
\newcommand{\cvprojectgap}{\par\addvspace{4pt}}

% \cvcert{name}{issuer}{meta line}
\newcommand{\cvcert}[3]{%
  \par\addvspace{3pt}%
  \needspace{3\baselineskip}%  name + issuer + meta line stay on one page
  \noindent
  \begin{tabularx}{\textwidth}{@{}Y r@{}}
    {\bfseries\color{cvaccent}#1} & {\footnotesize\color{cvmuted}#2}\\
    \multicolumn{2}{@{}l@{}}{\footnotesize\color{cvmuted}#3}
  \end{tabularx}}

% \cvscore{label}{score} - languages / test scores.
\newcommand{\cvscore}[2]{%
  \par\addvspace{2.5pt}\noindent
  \begin{tabularx}{\textwidth}{@{}Y r@{}}
    #1 & {\bfseries\color{cvaccent}#2}
  \end{tabularx}}

% \cvlabel{Responsibilities} - small bold lead-in.
\newcommand{\cvlabel}[1]{\par\addvspace{3pt}\noindent\textbf{#1}\par\addvspace{0pt}}

% \cvlevel{Advanced} - skill proficiency chip (small caps, accent colour).
\newcommand{\cvlevel}[1]{{\footnotesize\color{cvaccent}\textsc{#1}}}

% \cvtech{...} - shaded "Technologies & tools" strip.
\newcommand{\cvtech}[1]{%
  \par\addvspace{2.5pt}%
  {\setlength{\fboxsep}{4pt}%
   \noindent\colorbox{cvtechbg}{%
     \parbox{\dimexpr\textwidth-2\fboxsep\relax}{\textbf{Technologies \& tools:} #1}}}%
  \par\addvspace{0pt}}

% Skills table columns: bold category | items (flexible) | right-aligned level.
% They are plain column types (tabularx cannot be wrapped in \newenvironment).
\newcolumntype{K}{>{\bfseries\raggedright\arraybackslash}p{0.26\textwidth}}
% NOTE: an `m` (middle-aligned) column keeps the level on the text baseline;
% a `p{...}` column with \raggedleft would sit ~0.7\baselineskip too low.
\newcolumntype{S}{>{\raggedleft\arraybackslash\footnotesize\color{cvaccent}\scshape}%
                   m{0.17\textwidth}}
% Skills table: must be a macro that takes the body as its argument - tabularx
% scans the input stream for \end{tabularx}, so it cannot live in \cvskillend.
\newcommand{\cvskilltable}[1]{%
  \noindent\begin{tabularx}{\textwidth}{@{}K Y S@{}}#1\end{tabularx}}

% Bullets: tight and predictable.
\setlist[itemize]{leftmargin=1.35em,topsep=1pt,itemsep=0.5pt,parsep=0pt}

% --------------------- Experience & Projects entries -------------------------
% One list holds roles, projects or both, so each entry gets a project-style
% headline plus the role/company subtitle underneath it:
%   \cvxpentry{name}{subtitle (role $\cdot$ company)}{dates}
\newcommand{\cvxpentry}[3]{%
  \par\addvspace{5pt}%
  \needspace{5\baselineskip}%  headline + subtitle + description stay together
  \noindent
  \begin{tabularx}{\textwidth}{@{}Y r@{}}
    {\large\bfseries\color{cvaccent}#1} & {\footnotesize\color{cvmuted}#3}\\
    \multicolumn{2}{@{}l@{}}{\itshape\color{cvmuted}#2}
  \end{tabularx}%
  \par\addvspace{0pt}}
\newcommand{\cvxpgap}{\par\addvspace{5pt}}

% \cvdesc{...} - description paragraph under a \cvxpentry.
\newcommand{\cvdesc}[1]{\par\addvspace{2.5pt}\noindent #1\par\addvspace{0pt}}

% ------------------------------- Footer -------------------------------------
% One-line provenance note printed at the bottom of the LAST page. The document
% uses \pagestyle{empty}, so the PDF carries no page number and this line is the
% only thing in the footer area. Call it at the very end of src/cv.tex.
\newcommand{\cvfooter}[1]{%
  \par\vfill
  \noindent\textcolor{cvrule}{\rule{\textwidth}{0.8pt}}\par\addvspace{3pt}%
  {\footnotesize\color{cvmuted}#1}\par}

